%================================================================
%CONTENTS: Construction of Q from Z, Ring Axioms, Equivalence, Well Definedness, Mult Inverse/Division
%================================================================
In order to construct rational polynominals, we first must formally define rationals.
\subsection{Construction}

Define $\bQ \subseteq \mathbb{Z}\times \mathbb{Z}$, where the elements of $\mathbb{Q}$ are of the form $(a,b)$ where $a,b\in\mathbb{Z}, b\neq 0$. Define two operations $+,\cdot$, such that 

\begin{itemize}
    \item $+:\mathbb{Q}\times \mathbb{Q}\rightarrow \mathbb{Q}$ where $(a,b)+(c,d)=(ad+bc,bd)$
    \item $\cdot:\mathbb{Q}\times \mathbb{Q}\rightarrow \mathbb{Q}$ where $(a,b)\cdot (c,d)=(ac,bd)$
\end{itemize}

Also, define an equivalence relation on $\mathbb{Q}$. Two elements $(a,b),(c,d)$ in $\mathbb{Q}$ are equal iff $ad=bc$. Denote equivalence as $(a,b)\sim(c,d)$.

We can also denote $(a,b)$ as $\frac{a}{b}$.
The rationals are a Ring as defined by the ring axioms.
\subsection{Ring Axioms}
\textbf{Commutativity of Multiplication: } \[(a,b) \cdot (c,d) = (ac,bd) \qquad \text{and} \qquad (c,d) \cdot (a,b) = (ca, db) = (ac,bd)\]

\noindent\textbf{Commutativity of Addition: } \[(a,b)+(c,d) = (ad+bc, bd) \qquad \text{and} \qquad (c,d) + (a,b) = (cb+ad, db) = (ad+bc, bd)\]

\noindent\textbf{Associativity of Addition: } \[
((a,b) + (c,d))+(e,f) = (a,b)+((c,d)+(e,f))
\]
For the left hand side, \[   (ad+bc, bd)+(e,f) = (adf+bcf+bde, bdf)\]
For the right hand side \[(a,b) + (cf+de, df) = (adf+bcf+bde, bdf).\]



\noindent\textbf{Associativity of Multiplication: }
 \[(a,b)\cdot((c,d) \cdot(e,f)) = ((a,b) \cdot (c,d))\cdot(e,f)\]
\noindent
For the left hand side, \[(a,b)\cdot(ce,df)=(ace,bdf)\]
For the right hand side, 
\[
(ac,bd)\cdot(e,f)=(ace,bdf).
\]

\noindent\textbf{Additive Identity: }
\[
(a,b)+(0,1) = (a,b)
\]
From the definition, $(a,b)+(0,1)$ equals 
\[
(a\cdot 1+b\cdot 0, b\cdot 1)
\]
By the one axiom and zero axiom,
\[
(a,b)
\]
\noindent\textbf{Multiplicative Identity: }

\[
(a,b)\cdot(1,1) = (a,b)
\]
From the definition, $$(a,b)\cdot(1,1)$$ equals
\[
(a\cdot1,b\cdot1)
\]
By the one axiom, 
\[(a,b)\]
 
\noindent\textbf{Distributivity: }

We prove for $x,y,z\in \bQ$, we have $x(y+z)=xy+xz$. Suppose $x=(a,b),y=(c,d),z=(e,f)$. Then $x(y+z)=(a,b)\cdot((c,d)+(e,f))=(a,b)\cdot(cf+de,df)=(a(cf+de),b(df))=(acf+ade,bdf)$. We also have $xy=(a,b)\cdot(c,d)=(ac,bd)$ and $xz=(a,b)\cdot(e,f)=(ae,bf)$. We then have $xy+xz=(ac,bd)+(ae,bf)=(acbf+bdae,bdbf)=(acf+dae,dbf)=(acf+ade,bdf)$. Hence $x(y+z)=xy+xz$. Right distributivity follows immediately, since $(y+z)x=x(y+z)=xy+xz=yx+zx$.

\subsection{Equivalence}
Two rationals are equal if and only if they can be simplified to each other. In other words, 
\[
(a,b)\sim(c,d) \longleftrightarrow ad=bc
\]
\noindent\textbf{Symmetry: } If $(a,b)\sim(c,d)$, then $(c,d)\sim(a,b)$
\begin{proof}
    $ad=bc$. This can be written as $cb=da$, so 
\[
(c,d)\sim(a,b)
\]
\end{proof}
\noindent\textbf{Transitivity: } If $(a,b)\sim(c,d),(c,d)\sim(e,f)$ then $(a,b)\sim(e,f)$.
\begin{proof}
    Since $ad=bc, cf=de$, $adcf=bcde$ and $cdaf=cdbe$. Thus, $af=be$ so $(a,b)\sim(e,f)$
\end{proof}

\subsection{Well-Defined}
To show that addition is well defined, we show that $(a,b) + (c,d) = (a,b)+(e,f)$ if $(c,d) \sim (e,f)$. According to our criterion for an equivalence relation, $cf = de$. Therefore, we want to show $(ad+bc, bd) = (af+be, bf)$. To show these are equivalent, we want to show $abdf+b^2ed=abdf+b^2cf$. We cancel $abdf$ and $b^2$ because $b$ is nonzero to show that $de=cf$, which is true by our equivalence relation. \\

To show that multiplication is well defined, we show that $(a,b) \cdot (c,d) = (a,b)\cdot(e,f)$ if $(c,d) \sim (e,f)$. According to our criterion for an equivalence relation, $cf = de$. Therefore, we want to show $(ac, bd) = (ae, bf)$. To show these are equivalent, we want to show $acbf=bdae$. Suppose $a=0$. Then $0=0$, which is true. Otherwise, suppose $a\neq0$. Then because $ab=ac\implies b=c$ when $a\neq0$, we can cancel $ab$ from both sides to get $cf=de$, which is true by the equivalence $(c,d)=(e,f)$. Hence multiplication is well defined.




\subsection{Multiplicative Inverse}
%% whats the theorem
For any $(a,b)$ and $(c,d)\in\mathbb{Q}$, and $a\neq0$, $\exists(bc,ad)$ where
\[
(a,b)(bc,ad)=(abc,abd)
\]
and
\[
(abc,abd)\sim(c,d)
\]
then
\[
(a,b)\mid(c,d)
\]
Therefore, define $\frac{p}{q}$ for $p,q\in\mathbb{Q}$ and $p\neq0$ to be the unique solution $x$ to $px=q$. Note also that this is equivalent to the fraction notation of rationals since $\frac{(a,1)}{(b,1)}=(a,b)=\frac{a}{b}$. Then, define $p^{-1}$ to be $\frac{(1,1)}{p}$.
